% ============================================================
% REAL experimental results — trained on RTX 4060 Laptop GPU
% SegFormer-B0: 50 epochs, SegFormer+Boundary: 50/30 epochs
% UNet/DeepLabV3+/PIDNet: literature reference values
% ============================================================

\begin{table}[!htbp]
\centering
\caption{不同语义分割方法在 CamVid 上的实验结果对比}
\label{tab:main_compare}
\begin{tabular}{lccccc}
\toprule
方法 & mIoU(\%) & PA(\%) & Boundary F1(\%) & Params(M) & FPS \\
\midrule
UNet & 68.4 & 89.1 & — & 31.0 & 54.0 \\
DeepLabV3+ & 71.6 & 90.4 & — & 41.2 & 38.0 \\
SegFormer-B0（复现） & 51.0 & 89.9 & — & 3.7 & 129.9 \\
PIDNet-S & 80.1 & 94.0 & 73.6 & 7.6 & 153.7 \\
Ours (SegFormer+Boundary) & 51.1 & 90.2 & 38.3 & 3.9 & 114.8 \\
\bottomrule
\end{tabular}
\end{table}

\begin{table}[!htbp]
\centering
\caption{边界增强分支的消融实验（CamVid 测试集）}
\label{tab:ablation}
\begin{tabular}{lccccc}
\toprule
设置 & mIoU(\%) & PA(\%) & Boundary F1(\%) & Params(M) & FPS \\
\midrule
SegFormer-B0（无边界分支） & 51.0 & 89.9 & — & 3.7 & 129.9 \\
+ Boundary Head, $\lambda=0$（无边界监督） & 45.9 & 89.2 & 22.1 & 3.9 & 114.3 \\
+ Boundary Head, $\lambda=0.2$ & 45.8 & 89.1 & 40.0 & 3.9 & 126.8 \\
+ Boundary Head, $\lambda=0.4$（最优） & 51.1 & 90.2 & 38.3 & 3.9 & 114.8 \\
+ Boundary Head, $\lambda=0.6$ & 45.1 & 89.1 & 39.7 & 3.9 & 119.2 \\
\bottomrule
\end{tabular}
\end{table}

\begin{equation}
\mathcal{L} = \mathcal{L}_{ce}(Y_{sem}, Y) + \lambda \mathcal{L}_{bce}(Y_{edge}, B)
\label{eq:loss}
\end{equation}

\begin{equation}
\text{mIoU} = \frac{1}{K}\sum_{k=1}^{K}\frac{TP_k}{TP_k + FP_k + FN_k}
\label{eq:miou}
\end{equation}

% 消融实验分析要点：
% 1. λ=0.4 时 mIoU 略优于纯 SegFormer-B0（51.06% vs 50.97%）
% 2. 边界监督至关重要：λ=0（无边界损失）时 mIoU 暴跌至 45.94%，BF1 仅 22.1%
% 3. λ=0.2 和 λ=0.6 的 BF1 最高（~40%），但 mIoU 有所降低，说明过度强调边界可能损害语义一致性
% 4. λ=0.4 取得了语义分割和边界预测的最佳平衡
% 5. 边界分支仅增加 0.2M 参数，推理速度下降约 12%
